% Capítulo 3 (Principais Trabalhos Relacionados) — Slide 1 — Da literatura à lacuna
\begin{frame}{Da literatura à lacuna}
  \tiny
  \setlength{\leftmargini}{1.0em}
  \begin{columns}[t,onlytextwidth]
    \begin{column}{0.32\linewidth}
      \tabstyle
      \begin{beamercolorbox}[wd=\linewidth,sep=0.35em,rounded=true]{block title}
        \scriptsize\textbf{Políticas clássicas}
      \end{beamercolorbox}
      \vspace{0.3em}
      \textbf{Resolve:}
      \begin{itemize}\setlength\itemsep{0pt}\setlength\parskip{0pt}
        \item regras conhecidas de reposição
        \item interpretação simples
        \item baixo custo computacional
      \end{itemize}
      \textbf{Falta:}
      \begin{itemize}\setlength\itemsep{0pt}\setlength\parskip{0pt}
        \item adaptação a séries heterogêneas
        \item escolha contextual entre políticas
      \end{itemize}
    \end{column}
    \begin{column}{0.32\linewidth}
      \tabstyle
      \begin{beamercolorbox}[wd=\linewidth,sep=0.35em,rounded=true]{block title}
        \scriptsize\textbf{Previsão e otimização}
      \end{beamercolorbox}
      \vspace{0.3em}
      \textbf{Resolve:}
      \begin{itemize}\setlength\itemsep{0pt}\setlength\parskip{0pt}
        \item melhora estimativas de demanda
        \item calibra parâmetros de uma política já escolhida
      \end{itemize}
      \textbf{Falta:}
      \begin{itemize}\setlength\itemsep{0pt}\setlength\parskip{0pt}
        \item escolher qual família de política aplicar
        \item integrar custo, serviço e estabilidade em um processo de seleção
      \end{itemize}
    \end{column}
    \begin{column}{0.32\linewidth}
      \tabstyle
      \begin{beamercolorbox}[wd=\linewidth,sep=0.35em,rounded=true]{block title}
        \scriptsize\textbf{RL e métodos aprendidos}
      \end{beamercolorbox}
      \vspace{0.3em}
      \textbf{Resolve:}
      \begin{itemize}\setlength\itemsep{0pt}\setlength\parskip{0pt}
        \item aprende decisões sequenciais em ambientes simulados
        \item pode capturar dinâmicas complexas
      \end{itemize}
      \textbf{Falta:}
      \begin{itemize}\setlength\itemsep{0pt}\setlength\parskip{0pt}
        \item exige histórico/interações longas
        \item pouco validado em varejo real intermitente
        \item não resolve sozinho a seleção entre famílias
      \end{itemize}
    \end{column}
  \end{columns}
  \vspace{0.5em}
  \keymessage{A lacuna não é a ausência de políticas; é a ausência de um mecanismo contextual para escolher qual política aplicar em cada série loja-produto.}
  \note{
    A literatura oferece blocos sólidos para cada parte do problema, mas
    isolados. As políticas clássicas são simples e interpretáveis, mas se
    adaptam mal a séries heterogêneas. A previsão de demanda e a
    parametrização por meta-heurísticas melhoram estimativas e calibram uma
    política já escolhida, mas nenhuma das duas decide qual família de
    política aplicar. O aprendizado por reforço aprende decisões
    sequenciais, mas exige históricos longos e ainda é pouco validado em
    varejo real intermitente. Nenhum desses blocos resolve, isoladamente, a
    escolha contextual entre famílias heterogêneas de política -- e é
    exatamente nesse espaço que as três lacunas a seguir se tornam visíveis.
  }
\end{frame}
